\makeabstract
{
ROR1 CAR T-Cells As a Novel Treatment for Prostate Cancer
}
{
Sophia Bredar (1,2), Aerielle Matsangos, MHS (3), Djordje Atanackovic, MD (3)
}
{
\textsuperscript{1}Amherst College, Amherst, MA\\
\textsuperscript{2}Nathan Schnaper Intern Program in Translational Cancer Research, University of Maryland, Baltimore, MD, USA\\
\textsuperscript{3}University of Maryland Marlene \& Stewart Greenebaum Comprehensive Cancer Center, Baltimore, MD, USA
}
{
Chimeric Antigen Receptor (CAR) T-cells engineered to kill tumor cells have revolutionized the treatment of hematologic cancers. To do this, CAR T-cells use an artificially expressed antibody to bind to surface antigen on tumors. The efficacy of this approach against solid tumors has been limited, which is primarily due to (1) the absence of appropriate targets from solid tumors and (2) the presence of immunosuppressive factors like TGFβ in the tumor microenvironment. Our lab is for the first time investigating the tumor-targeting ability of novel CAR T-cells specific for the Receptor Tyrosine Kinase-like Orphan Receptor 1 (ROR1) antigen.
}
{
We analyzed ROR1 surface expression on various prostate cancer cell lines by flow cytometry. Next, tumor cells underwent GFP-luc lentivirus transduction and were sorted to yield pure target cells for cytotoxicity assays. Finally, we evaluated the tumor-killing ability of ROR1-specific CAR T-cells and compared them to ROR1 CAR T-cells armored with a dominant-negative form of TGFβRII.
}
{
Our novel CAR T-cells were able to bind to and produce various cytokines in response to recombinant ROR1 protein. We found that most of the prostate tumor cell lines express substantial ROR1 levels. Finally, unmodified and “armored” CAR T-cells both efficiently killed prostate cancer cells, even at low effector/target ratios.
}
{
Our findings suggest for the first time that ROR1 CAR T-cells may serve as a promising treatment for prostate cancer, especially when combined with anti-TGFβ mechanisms. Additional studies are underway in our lab evaluating this novel CAR T-cell approach in other solid tumors.
}

\newpage
